\chapter{Mitral Valve Repair or Replacement}
\section*{What Is This Test or Procedure?}
Mitral valve surgery repairs a diseased valve or replaces it with a mechanical or tissue prosthesis.
\section*{Why This Test or Procedure Is Ordered}
It may be required for severe symptomatic mitral regurgitation, stenosis, infection, or progressive effects on heart chambers.
\section*{How This Test or Procedure Is Performed}
Open or minimally invasive surgery is performed under cardiopulmonary bypass. Repair may reshape leaflets or support structures; replacement removes the native valve and implants a prosthesis.
\section*{Risks and Follow-up}
Bleeding, stroke, infection, arrhythmia, kidney injury, residual leakage, recurrent disease, and prosthetic complications are possible. Anticoagulation may be required for a mechanical valve.
\section*{References}
\begin{enumerate}\item MedlinePlus. Heart Valve Surgery. U.S. National Library of Medicine; 2025.\item Current Medical Diagnosis and Treatment. McGraw-Hill; 2025.\end{enumerate}
\clearpage
